% ─────────────────────────────────────────────────────────────────────────────
% Workload–Recommender Alignment: A Theoretical Framework
% To include in the main paper: % ─────────────────────────────────────────────────────────────────────────────
% Workload–Recommender Alignment: A Theoretical Framework
% To include in the main paper: % ─────────────────────────────────────────────────────────────────────────────
% Workload–Recommender Alignment: A Theoretical Framework
% To include in the main paper: % ─────────────────────────────────────────────────────────────────────────────
% Workload–Recommender Alignment: A Theoretical Framework
% To include in the main paper: \input{workload_theory}
% Place between the framework description and the Research Challenges section.
% ─────────────────────────────────────────────────────────────────────────────

\section{Theoretical Analysis: Workload--Recommender Alignment}
\label{sec:theory}

The purpose of this section is to give the benchmark results a principled
foundation: we first construct a formal
model of EDA sessions and define precisely what each recommender optimises, then
derive conditions on the \emph{workload} under which each recommender succeeds
or fails.  The analysis yields both an explanation of the benchmark design and a
prescriptive recipe for constructing controlled evaluations targeted at any
recommender.

%─────────────────────────────────────────
\subsection{Formal Session Model}
\label{sec:theory:model}

\paragraph{Schema and tuples.}
Let $\mathcal{A} = \{A_1, \ldots, A_d\}$ be a fixed relational schema.  Each
attribute $A_i$ has a finite domain $\text{dom}(A_i)$.  A \emph{tuple}
$t \in \mathcal{T}$ is a complete assignment $t : \mathcal{A} \to
\bigcup_i \text{dom}(A_i)$ such that $t(A_i) \in \text{dom}(A_i)$.  The
\emph{fact table} $\mathcal{T}$ is a finite multiset of such tuples.

\paragraph{Association rules.}
An \emph{association rule} $r = (X \Rightarrow Y)$ is a pair of attribute–value
sets $X, Y \subset \{A_i = v \mid A_i \in \mathcal{A},\, v \in \text{dom}(A_i)\}$
with $X \cap Y = \emptyset$.  A tuple $t$ \emph{satisfies the antecedent} $X$
if $t(A_i) = v$ for all $(A_i = v) \in X$, and similarly for $Y$.  Standard
interestingness measures are:
\begin{align}
  \text{supp}(r) &= \frac{|\{t \in \mathcal{T} \mid t \models X \cup Y\}|}{|\mathcal{T}|},
  \label{eq:support}\\
  \text{conf}(r) &= \frac{\text{supp}(r)}{\text{supp}(X)},
  \label{eq:confidence}\\
  \text{lift}(r) &= \frac{\text{conf}(r)}{\text{supp}(Y)},
  \label{eq:lift}
\end{align}
where $\text{supp}(X)$ (resp. $\text{supp}(Y)$) denotes the fraction of tuples
satisfying $X$ (resp. $Y$) alone.  The \emph{J-measure}~\cite{smyth1992rule}
is a mutual-information proxy:
\begin{equation}
  J(r) = \text{supp}(r)\log\frac{\text{conf}(r)}{\text{supp}(Y)}
        + (1-\text{supp}(r))\log\frac{1-\text{conf}(r)}{1-\text{supp}(Y)}.
  \label{eq:jmeasure}
\end{equation}

\paragraph{EDA session and prediction task.}
An \emph{EDA session} of length $Q$ is an ordered sequence of query results
$\mathcal{S} = (R_0, R_1, \ldots, R_{Q-1})$ where each $R_q \subseteq \mathcal{T}$
is a finite multiset.  The \emph{history} at step $q$ is
$\mathcal{H}_q = (R_0, \ldots, R_{q-1})$.

At each step $q$ with \emph{prediction gap} $g \geq 1$, a recommender is given
$(R_q, \mathcal{H}_q)$ and must rank the tuples of $R_q$.  The ground-truth
relevance label for a tuple $t \in R_q$ is
\begin{equation}
  \text{rel}(t,q,g) = \mathbf{1}\bigl[t \in R_{q+g}\bigr].
  \label{eq:relevance}
\end{equation}
That is, a tuple is relevant if it \emph{reappears} in the result set $g$ steps
into the future --- a natural formalisation of ``the analyst will want to see
this again soon''.

%─────────────────────────────────────────
\subsection{Evaluation: nDCG as the Quality Criterion}
\label{sec:theory:ndcg}

Let $\pi$ be the ranking of $R_q$ produced by a recommender and let $k$ be the
recommendation budget.  The normalised discounted cumulative gain at rank $k$
is
\begin{equation}
  \text{nDCG}_k(\pi, q, g)
  = \frac{1}{Z_k}
    \sum_{i=1}^{k}
    \frac{\text{rel}(\pi(i), q, g)}{\log_2(i+1)},
  \label{eq:ndcg}
\end{equation}
where $Z_k$ is the ideal DCG (all relevant tuples at the top).  We write
$\text{nDCG}_k(B, \mathcal{S})$ for the mean over all valid $(q, g)$ pairs in
session $\mathcal{S}$ for recommender $B$.

\paragraph{Overlap concentration.}
Define the \emph{rule-conforming overlap fraction} of a query pair $(q, q+g)$
as
\begin{equation}
  \rho(q,g)
  = \frac{|R_q \cap R_{q+g} \cap \mathcal{T}^*|}{|R_q \cap R_{q+g}|},
  \label{eq:rho}
\end{equation}
where $\mathcal{T}^* = \{t \in \mathcal{T} \mid \exists r : t \models X_r
\cup Y_r\}$ is the set of rule-conforming tuples.  When $\rho(q,g) \approx 1$,
predicting overlap is essentially the same as identifying rule tuples inside
$R_q$.

%─────────────────────────────────────────
\subsection{Recommender Scoring Functions}
\label{sec:theory:scorers}

We model each recommender as a \emph{scoring function}
$s_B : R_q \times \mathcal{H}_q \to \mathbb{R}$ that induces a ranking $\pi_B$
over $R_q$.

\paragraph{Random.}
$s_{\mathrm{rand}}(t) \sim \mathrm{Uniform}(0,1)$, independent of everything.

\paragraph{Frequency.}
Let $\text{freq}_q(A_i = v) = |\{(s, u) \mid s < q,\, u \in R_s,\, u(A_i) = v\}|$
be the historical count of value $v$ in attribute $A_i$ across past results.
\begin{equation}
  s_{\mathrm{freq}}(t) = \frac{1}{d}\sum_{i=1}^d \text{freq}_q\bigl(A_i = t(A_i)\bigr).
  \label{eq:freq-score}
\end{equation}

\paragraph{Clustering.}
One-hot encode each tuple in $R_q$ to a vector $\phi(t) \in \{0,1\}^D$.  Run
$K$-means on $\{\phi(t) \mid t \in R_q\}$ to obtain centroids
$\{\mu_1, \ldots, \mu_K\}$.  Score each tuple by proximity to its assigned
centroid's \emph{neighbourhood density}:
\begin{equation}
  s_{\mathrm{clust}}(t) = -\|\phi(t) - \mu_{k(t)}\|^2,
  \label{eq:clust-score}
\end{equation}
where $k(t) = \arg\min_k \|\phi(t) - \mu_k\|$.  Tuples near a centroid score
high; outliers score low.

\paragraph{Similarity.}
Let $\bar{\phi}_q = \frac{1}{|R_q|}\sum_{t \in R_q}\phi(t)$ be the mean
encoding of the current result.  Rank by cosine similarity to this centroid:
\begin{equation}
  s_{\mathrm{sim}}(t) = \frac{\phi(t) \cdot \bar{\phi}_q}{\|\phi(t)\|\,\|\bar{\phi}_q\|}.
  \label{eq:sim-score}
\end{equation}

\paragraph{Sampling.}
A stratified random sample of $k$ tuples from $R_q \cup \mathcal{H}_{q}^{\rm flat}$
(the union of current and past results), with each tuple weighted by its
historical frequency.  This interpolates between Random and Frequency depending
on the history length.

\paragraph{MDI.}
The MDI score accumulates three components across the history:
\begin{equation}
  s_{\mathrm{MDI}}(t, q)
  = \alpha \cdot A(t, q) + \beta \cdot D(t, q) + \gamma \cdot N(t, q),
  \label{eq:mdi-score}
\end{equation}
where $\alpha + \beta + \gamma = 1$.

\begin{description}
  \item[Association component $A(t,q)$:]
    Let $\mathcal{R}_q$ be association rules mined from the history up to step
    $q$.  For rule $r \in \mathcal{R}_q$, let
    $w_r(q) = e^{-\lambda_r (q - q_r)}$ be an exponential decay weight
    ($\lambda_r$ is the decay rate, $q_r$ the query at which $r$ was mined).
    The composite rule score is
    \begin{equation}
      I_{\mathrm{assoc}}(r)
      = \lambda_1\,\mathrm{conf}(r) + \lambda_2\,\mathrm{supp}(r)
      + \lambda_3\,\mathrm{lift}(r) + \lambda_4\,J(r),
      \label{eq:Iassoc}
    \end{equation}
    and the association component for tuple $t$ is
    \begin{equation}
      A(t,q)
      = \sum_{r \in \mathcal{R}_q}
        w_r(q)\,\frac{I_{\mathrm{assoc}}(r)}{|X_r \cup Y_r|}\,
        \mathbf{1}[t \models X_r \cup Y_r].
      \label{eq:Acomponent}
    \end{equation}

  \item[Diversity component $D(t,q)$:]
    Measures how much adding $t$ to a candidate summary would increase the
    entropy of attribute-value distributions in $\mathcal{H}_q$ (combining
    Shannon, Simpson, Gini, Berger--Parker, and McIntosh indices with weights
    $\mu_1, \ldots, \mu_5$).

  \item[Novelty component $N(t,q)$:]
    Penalises tuples whose attribute values are over-represented in
    $\mathcal{H}_q$, favouring attribute–value combinations that have appeared
    sparsely in the session history.
\end{description}

%─────────────────────────────────────────
\subsection{Four Workload Parameters}
\label{sec:theory:params}

The behaviour of each recommender is controlled by four workload parameters
that can be tuned independently in a synthetic benchmark:

\begin{description}
  \item[$\rho$ --- signal fraction:]
    The fraction of tuples in typical large query results that are
    rule-conforming:
    $\rho = \mathbb{E}[|R_q \cap \mathcal{T}^*| / |R_q|]$.
    In our benchmark, $\rho \approx 0.20$--$0.30$ for large queries.

  \item[$\xi$ --- frequency skew:]
    The ratio of the empirical count of common values to the count of rare
    (rule-antecedent) values in $\mathcal{H}_q$:
    $\xi = \text{freq}_q(\text{common}) / \text{freq}_q(\text{rare})$.
    Controlled by the bias probability $p_{\mathrm{common}}$ (set to $0.80$).

  \item[$\sigma^2_{\mathrm{noise}}$ --- noise cluster variance:]
    The mean intra-cluster variance of noise tuples in encoded feature space.
    When $p_{\mathrm{common}}$ is high, the noise distribution concentrates on
    $2$--$3$ values per attribute, making $\sigma^2_{\mathrm{noise}}$ small.

  \item[$Q^*$ --- required memory depth:]
    The minimum session length (number of queries before the evaluation gap)
    needed for MDI to mine association rules with sufficient frequency to
    produce reliable confidence estimates.  Depends on $\text{supp}(r)$ and
    the minimum support threshold $\theta_s$.
\end{description}

%─────────────────────────────────────────
\subsection{Failure Conditions for Baseline Recommenders}
\label{sec:theory:failures}

We now state conditions under which each baseline fails to identify
rule-conforming overlap tuples, giving MDI a ranking advantage.

\begin{proposition}[Frequency failure]
\label{prop:freq}
Let $f_c = \text{freq}_q(\text{common})$ and $f_r = \text{freq}_q(\text{rare})$
be the mean historical counts of common and rare attribute values respectively.
Fix a query result $R_q$ containing noise tuples $R_q^{\rm noise}$ biased toward
common values and rule tuples $R_q^* \subseteq \mathcal{T}^*$.
If $\xi = f_c / f_r > 1$, then for every rule tuple $t^* \in R_q^*$ and typical
noise tuple $t^n \in R_q^{\rm noise}$,
\[
  s_{\mathrm{freq}}(t^n) > s_{\mathrm{freq}}(t^*).
\]
\textit{Proof sketch.}
By construction, each component $\text{freq}_q(A_i = t^n(A_i)) = f_c$ (biased
toward common values) while $\text{freq}_q(A_i = t^*(A_i)) = f_r$ for the
rule-specific attributes.  The inequality follows directly from Eq.~\eqref{eq:freq-score}
and $\xi > 1$.
\end{proposition}

\begin{proposition}[Clustering failure]
\label{prop:clust}
Suppose the noise distribution concentrates on $c$ values per attribute
($c \ll |\text{dom}(A_i)|$) so that $\sigma^2_{\mathrm{noise}} \leq \delta$
for a small $\delta > 0$, and let rule tuples scatter uniformly across
$|\text{dom}(A_i)|$ values.  Then for any $K$-means solution with $K$ clusters,
at least $K - O(\sqrt{\delta})$ centroids coincide with noise cluster means, and
every rule tuple $t^*$ receives an assignment distance
$\|{\phi}(t^*) - \mu_{k(t^*)}\|^2 \geq \Omega(1/c)$, while
$\|{\phi}(t^n) - \mu_{k(t^n)}\|^2 \leq \delta$ for noise tuples.
By Eq.~\eqref{eq:clust-score}, $s_{\mathrm{clust}}(t^n) > s_{\mathrm{clust}}(t^*)$.
\end{proposition}

\begin{proposition}[Similarity failure]
\label{prop:sim}
When $|R_q^{\rm noise}| \gg |R_q^*|$, the result centroid
$\bar{\phi}_q$ is dominated by noise encodings.  Since rule tuples use
attribute values orthogonal to common values in the one-hot representation,
$\phi(t^*) \cdot \bar{\phi}_q \approx \rho\,\|\phi(t^*)\|$, whereas
$\phi(t^n) \cdot \bar{\phi}_q \approx (1-\rho)\,\|\phi(t^n)\|$.  For
$\rho < 1/2$, the cosine score satisfies $s_{\mathrm{sim}}(t^n) > s_{\mathrm{sim}}(t^*)$.
\end{proposition}

\begin{proposition}[Random failure]
\label{prop:rand}
An unbiased random ranking of $n = |R_q|$ tuples, of which $m = |R_q^*|$ are
relevant (ground truth), achieves expected nDCG
\[
  \mathbb{E}[\text{nDCG}_k] \leq \frac{m/n \cdot H_k}{Z_k},
\]
where $H_k = \sum_{i=1}^k 1/\log_2(i+1)$ is the un-normalised ideal DCG.
This equals the ideal nDCG only when $m/n = 1$; for $\rho = m/n \ll 1$
(sparse ground truth), random performance degrades proportionally to $\rho$.
\end{proposition}

\paragraph{Remark on Sampling.}
Sampling interpolates between Random and Frequency: with sufficient history it
converges to the Frequency scorer.  Therefore Propositions~\ref{prop:freq}
and~\ref{prop:rand} bound its performance from above and below respectively,
and it fails for the same structural reason as Frequency once $\xi > 1$.

%─────────────────────────────────────────
\subsection{MDI's Temporal Accumulation Advantage}
\label{sec:theory:mdi}

The key distinction of MDI is that its association component $A(t,q)$
(Eq.~\ref{eq:Acomponent}) \emph{grows} with session history while the baselines
remain stationary or degrade.

\begin{proposition}[Temporal accumulation]
\label{prop:temporal}
Let $r^* = (X^* \Rightarrow Y^*)$ be a ground-truth embedded rule with
confidence $c^* = \text{conf}(r^*)$ and support $s^* = \text{supp}(r^*)$.
Suppose each large query in the session contains a \emph{persistent core} of
$m_c$ rule-conforming tuples drawn i.i.d. from $\mathcal{T}^*$, and MDI mines
$r^*$ from accumulated data at step $q$ with estimated confidence
$\hat{c}_q(r^*) \to c^*$ as $q \to \infty$.  Then the association component
for any tuple $t^* \models X^* \cup Y^*$ satisfies
\[
  A(t^*, q) \xrightarrow{q \to \infty}
  \frac{c^*\,(\lambda_1 + \lambda_3 \cdot c^*/s_Y^* + \lambda_4 \cdot J^*)}{|X^* \cup Y^*|},
\]
where $s_Y^* = \text{supp}(Y^*)$ and $J^*$ is the J-measure of $r^*$.
Moreover, since the rule decay weight $w_{r^*}(q) = e^{-\lambda_r(q - q_{r^*})}$
is refreshed every time $r^*$ is re-mined (i.e., at every large query), the
effective decay does not erode the signal over the session.
\end{proposition}

\begin{corollary}[Gap advantage]
\label{cor:gap}
Because $A(t^*, q)$ accumulates across queries regardless of the prediction gap
$g$, while $s_{\mathrm{freq}}(t)$, $s_{\mathrm{clust}}(t)$, and
$s_{\mathrm{sim}}(t)$ are memoryless (computed from $R_q$ and recent history
only), MDI's advantage over these baselines is \emph{non-decreasing} in session
length $q$ and \emph{independent} of $g$.  This means that at large gaps
($g = 5, 10$), MDI retains its advantage while memoryless baselines degrade
because the current result $R_{q+g}$ drifts from $R_q$ in ways that cannot be
predicted from $R_q$ alone.
\end{corollary}

%─────────────────────────────────────────
\subsection{Benchmark Construction: A Prescriptive Recipe}
\label{sec:theory:recipe}

The preceding propositions jointly characterise when MDI wins, but the
framework also prescribes how to \emph{deliberately construct workloads that
favour any given recommender}.  We state this as a theorem and summarise the
parameter regimes in Table~\ref{tab:recipe}.

\begin{theorem}[Workload-determined advantage]
\label{thm:recipe}
Let $\mathcal{B} = \{\text{MDI, Frequency, Clustering, Similarity, Random,
Sampling}\}$ be the recommender set.  For each $B \in \mathcal{B}$ there
exists a workload characterised by parameters $(\rho, \xi, \sigma^2_{\rm noise},
Q^*)$ such that $\text{nDCG}_k(B, \mathcal{S}) > \text{nDCG}_k(B', \mathcal{S})$
for all $B' \neq B$.  The dominating regime for each recommender is:
\begin{itemize}
  \item \textbf{MDI:} $\rho < 1/2$, $\xi \gg 1$,
        $\sigma^2_{\rm noise} \ll 1$, $Q \geq Q^*(s^*, \theta_s)$.
        Rule-conforming overlap tuples are rare, noise is homogeneous, and
        the session is long enough for rule accumulation.

  \item \textbf{Frequency:} $\rho \approx 1$ and $\xi < 1$ (rare values carry the
        overlap signal, but all overlap tuples have above-average frequency).

  \item \textbf{Clustering:} Overlap tuples form a tight, distinct cluster
        ($\sigma^2_{\rm overlap} \ll \sigma^2_{\rm noise}$) separated from noise
        in encoded space; $K \geq 2$.

  \item \textbf{Similarity:} The future result $R_{q+g}$ is \emph{more similar}
        to the current result centroid than any novel tuple; queries are
        monotonic refinements of the same analytical theme.

  \item \textbf{Random / Sampling:} $\rho \approx 1/2$ and the session is short
        ($Q < Q^*$); association rules are not yet reliable, so memory provides
        no signal and random selection is near-optimal.
\end{itemize}
\end{theorem}

\begin{table}[t!]
  \centering
  \caption{%
    Workload parameter regimes that favour each recommender.
    $\checkmark$ = condition required; $\times$ = not required
    (or reversed).
  }
  \label{tab:recipe}
  \footnotesize
  \setlength{\tabcolsep}{3pt}
  \begin{tabular}{l c c c c}
    \toprule
    \textbf{Rec.}
      & \makecell{$\rho \!\ll\! 1$\\[-2pt]\scriptsize sparse sig.}
      & \makecell{$\xi \!\gg\! 1$\\[-2pt]\scriptsize freq.\ skew}
      & \makecell{$\sigma^2_{\rm n} \!\ll\! 1$\\[-2pt]\scriptsize noise hom.}
      & \makecell{$Q \!\geq\! Q^*$\\[-2pt]\scriptsize deep hist.} \\
    \midrule
    MDI        & \checkmark & \checkmark & \checkmark & \checkmark \\
    Frequency  & $\times$   & $\times$   & $\times$   & \checkmark \\
    Clustering & $\times$   & $\times$   & $\times$\textsuperscript{a}  & $\times$ \\
    Similarity & $\times$   & $\times$   & $\times$\textsuperscript{b}  & $\times$ \\
    Random     & $\times$   & $\times$   & $\times$   & $\times$   \\
    Sampling   & $\times$   & $\times$   & $\times$   & $\times$   \\
    \bottomrule
    \multicolumn{5}{l}{\scriptsize \textsuperscript{a}distinct clusters \quad \textsuperscript{b}centroid-aligned}
  \end{tabular}
\end{table}

\paragraph{Instantiation in the benchmark generator.}
Table~\ref{tab:benchmark-params} shows how the parameters of
\texttt{generate\_benchmark.py} map to the theoretical quantities.  All four
conditions in Theorem~\ref{thm:recipe} for MDI are satisfied by construction:
$\rho \approx 0.22$ for large queries; $\xi = p_{\rm common}/(1 - p_{\rm common}) = 4$
($p_{\rm common} = 0.80$); noise rows concentrate on 2 values per attribute in a
9-attribute schema so $\sigma^2_{\rm noise} \leq 0.11$; and with $Q = 30$ queries
per session (each large query contributing $\sim 150$ tuples) the expected
antecedent support count after 10 large queries is
$10 \times 150 \times \text{supp}(r^*) \approx 90 > 1/\theta_s = 12.5$,
satisfying $Q \geq Q^*$.

\begin{table}[t!]
  \centering
  \caption{%
    Mapping from \texttt{generate\_benchmark.py} parameters to theoretical
    quantities.
  }
  \label{tab:benchmark-params}
  \footnotesize
  \setlength{\tabcolsep}{3pt}
  \begin{tabular}{l l l}
    \toprule
    \textbf{Theory param.} & \textbf{Generator param.} & \textbf{Value} \\
    \midrule
    $\rho$ (signal frac.)     & rule rows / lg.\ query sz.       & $\approx 0.22$ \\
    $\xi$ (freq.\ skew)       & $p_{\rm common} = 0.80$ bias     & $4.0$ \\
    $\sigma^2_{\rm noise}$    & 2 common vals / 9 attrs           & $\leq 0.11$ \\
    $\text{supp}(r^*)$        & 12\% rule rows; conf $= 0.85$    & $\approx 0.10$ \\
    $Q^*$                     & min\_supp $= 0.08$, 3000 rows    & $\approx 3$ lg.\ queries \\
    Core overlap $m_c$        & 20 persistent rule tuples         & $20$ \\
    Sm.\ query overlap        & 24--30 of 40--50 tuples           & $\approx 0.60$--$0.70$ \\
    \bottomrule
  \end{tabular}
\end{table}

\paragraph{Implications for real-world workloads.}
Real EDA sessions over operational databases frequently exhibit high $\xi$ simply
because data skew is ubiquitous, but they rarely satisfy $\rho \ll 1$ unless the
analyst is explicitly drilling into a rare sub-population.  This means MDI's
advantage is most pronounced in \emph{drill-down} scenarios --- exactly the kind
of session studied in the SDSS and SIMBA benchmarks --- and shrinks toward the
level of simpler baselines in broad exploratory sweeps where overlap is spread
uniformly across result rows.  The framework thus not only explains the benchmark
results but also signals where MDI deployment will deliver the greatest practical
benefit.

% Place between the framework description and the Research Challenges section.
% ─────────────────────────────────────────────────────────────────────────────

\section{Theoretical Analysis: Workload--Recommender Alignment}
\label{sec:theory}

The purpose of this section is to give the benchmark results a principled
foundation: we first construct a formal
model of EDA sessions and define precisely what each recommender optimises, then
derive conditions on the \emph{workload} under which each recommender succeeds
or fails.  The analysis yields both an explanation of the benchmark design and a
prescriptive recipe for constructing controlled evaluations targeted at any
recommender.

%─────────────────────────────────────────
\subsection{Formal Session Model}
\label{sec:theory:model}

\paragraph{Schema and tuples.}
Let $\mathcal{A} = \{A_1, \ldots, A_d\}$ be a fixed relational schema.  Each
attribute $A_i$ has a finite domain $\text{dom}(A_i)$.  A \emph{tuple}
$t \in \mathcal{T}$ is a complete assignment $t : \mathcal{A} \to
\bigcup_i \text{dom}(A_i)$ such that $t(A_i) \in \text{dom}(A_i)$.  The
\emph{fact table} $\mathcal{T}$ is a finite multiset of such tuples.

\paragraph{Association rules.}
An \emph{association rule} $r = (X \Rightarrow Y)$ is a pair of attribute–value
sets $X, Y \subset \{A_i = v \mid A_i \in \mathcal{A},\, v \in \text{dom}(A_i)\}$
with $X \cap Y = \emptyset$.  A tuple $t$ \emph{satisfies the antecedent} $X$
if $t(A_i) = v$ for all $(A_i = v) \in X$, and similarly for $Y$.  Standard
interestingness measures are:
\begin{align}
  \text{supp}(r) &= \frac{|\{t \in \mathcal{T} \mid t \models X \cup Y\}|}{|\mathcal{T}|},
  \label{eq:support}\\
  \text{conf}(r) &= \frac{\text{supp}(r)}{\text{supp}(X)},
  \label{eq:confidence}\\
  \text{lift}(r) &= \frac{\text{conf}(r)}{\text{supp}(Y)},
  \label{eq:lift}
\end{align}
where $\text{supp}(X)$ (resp. $\text{supp}(Y)$) denotes the fraction of tuples
satisfying $X$ (resp. $Y$) alone.  The \emph{J-measure}~\cite{smyth1992rule}
is a mutual-information proxy:
\begin{equation}
  J(r) = \text{supp}(r)\log\frac{\text{conf}(r)}{\text{supp}(Y)}
        + (1-\text{supp}(r))\log\frac{1-\text{conf}(r)}{1-\text{supp}(Y)}.
  \label{eq:jmeasure}
\end{equation}

\paragraph{EDA session and prediction task.}
An \emph{EDA session} of length $Q$ is an ordered sequence of query results
$\mathcal{S} = (R_0, R_1, \ldots, R_{Q-1})$ where each $R_q \subseteq \mathcal{T}$
is a finite multiset.  The \emph{history} at step $q$ is
$\mathcal{H}_q = (R_0, \ldots, R_{q-1})$.

At each step $q$ with \emph{prediction gap} $g \geq 1$, a recommender is given
$(R_q, \mathcal{H}_q)$ and must rank the tuples of $R_q$.  The ground-truth
relevance label for a tuple $t \in R_q$ is
\begin{equation}
  \text{rel}(t,q,g) = \mathbf{1}\bigl[t \in R_{q+g}\bigr].
  \label{eq:relevance}
\end{equation}
That is, a tuple is relevant if it \emph{reappears} in the result set $g$ steps
into the future --- a natural formalisation of ``the analyst will want to see
this again soon''.

%─────────────────────────────────────────
\subsection{Evaluation: nDCG as the Quality Criterion}
\label{sec:theory:ndcg}

Let $\pi$ be the ranking of $R_q$ produced by a recommender and let $k$ be the
recommendation budget.  The normalised discounted cumulative gain at rank $k$
is
\begin{equation}
  \text{nDCG}_k(\pi, q, g)
  = \frac{1}{Z_k}
    \sum_{i=1}^{k}
    \frac{\text{rel}(\pi(i), q, g)}{\log_2(i+1)},
  \label{eq:ndcg}
\end{equation}
where $Z_k$ is the ideal DCG (all relevant tuples at the top).  We write
$\text{nDCG}_k(B, \mathcal{S})$ for the mean over all valid $(q, g)$ pairs in
session $\mathcal{S}$ for recommender $B$.

\paragraph{Overlap concentration.}
Define the \emph{rule-conforming overlap fraction} of a query pair $(q, q+g)$
as
\begin{equation}
  \rho(q,g)
  = \frac{|R_q \cap R_{q+g} \cap \mathcal{T}^*|}{|R_q \cap R_{q+g}|},
  \label{eq:rho}
\end{equation}
where $\mathcal{T}^* = \{t \in \mathcal{T} \mid \exists r : t \models X_r
\cup Y_r\}$ is the set of rule-conforming tuples.  When $\rho(q,g) \approx 1$,
predicting overlap is essentially the same as identifying rule tuples inside
$R_q$.

%─────────────────────────────────────────
\subsection{Recommender Scoring Functions}
\label{sec:theory:scorers}

We model each recommender as a \emph{scoring function}
$s_B : R_q \times \mathcal{H}_q \to \mathbb{R}$ that induces a ranking $\pi_B$
over $R_q$.

\paragraph{Random.}
$s_{\mathrm{rand}}(t) \sim \mathrm{Uniform}(0,1)$, independent of everything.

\paragraph{Frequency.}
Let $\text{freq}_q(A_i = v) = |\{(s, u) \mid s < q,\, u \in R_s,\, u(A_i) = v\}|$
be the historical count of value $v$ in attribute $A_i$ across past results.
\begin{equation}
  s_{\mathrm{freq}}(t) = \frac{1}{d}\sum_{i=1}^d \text{freq}_q\bigl(A_i = t(A_i)\bigr).
  \label{eq:freq-score}
\end{equation}

\paragraph{Clustering.}
One-hot encode each tuple in $R_q$ to a vector $\phi(t) \in \{0,1\}^D$.  Run
$K$-means on $\{\phi(t) \mid t \in R_q\}$ to obtain centroids
$\{\mu_1, \ldots, \mu_K\}$.  Score each tuple by proximity to its assigned
centroid's \emph{neighbourhood density}:
\begin{equation}
  s_{\mathrm{clust}}(t) = -\|\phi(t) - \mu_{k(t)}\|^2,
  \label{eq:clust-score}
\end{equation}
where $k(t) = \arg\min_k \|\phi(t) - \mu_k\|$.  Tuples near a centroid score
high; outliers score low.

\paragraph{Similarity.}
Let $\bar{\phi}_q = \frac{1}{|R_q|}\sum_{t \in R_q}\phi(t)$ be the mean
encoding of the current result.  Rank by cosine similarity to this centroid:
\begin{equation}
  s_{\mathrm{sim}}(t) = \frac{\phi(t) \cdot \bar{\phi}_q}{\|\phi(t)\|\,\|\bar{\phi}_q\|}.
  \label{eq:sim-score}
\end{equation}

\paragraph{Sampling.}
A stratified random sample of $k$ tuples from $R_q \cup \mathcal{H}_{q}^{\rm flat}$
(the union of current and past results), with each tuple weighted by its
historical frequency.  This interpolates between Random and Frequency depending
on the history length.

\paragraph{MDI.}
The MDI score accumulates three components across the history:
\begin{equation}
  s_{\mathrm{MDI}}(t, q)
  = \alpha \cdot A(t, q) + \beta \cdot D(t, q) + \gamma \cdot N(t, q),
  \label{eq:mdi-score}
\end{equation}
where $\alpha + \beta + \gamma = 1$.

\begin{description}
  \item[Association component $A(t,q)$:]
    Let $\mathcal{R}_q$ be association rules mined from the history up to step
    $q$.  For rule $r \in \mathcal{R}_q$, let
    $w_r(q) = e^{-\lambda_r (q - q_r)}$ be an exponential decay weight
    ($\lambda_r$ is the decay rate, $q_r$ the query at which $r$ was mined).
    The composite rule score is
    \begin{equation}
      I_{\mathrm{assoc}}(r)
      = \lambda_1\,\mathrm{conf}(r) + \lambda_2\,\mathrm{supp}(r)
      + \lambda_3\,\mathrm{lift}(r) + \lambda_4\,J(r),
      \label{eq:Iassoc}
    \end{equation}
    and the association component for tuple $t$ is
    \begin{equation}
      A(t,q)
      = \sum_{r \in \mathcal{R}_q}
        w_r(q)\,\frac{I_{\mathrm{assoc}}(r)}{|X_r \cup Y_r|}\,
        \mathbf{1}[t \models X_r \cup Y_r].
      \label{eq:Acomponent}
    \end{equation}

  \item[Diversity component $D(t,q)$:]
    Measures how much adding $t$ to a candidate summary would increase the
    entropy of attribute-value distributions in $\mathcal{H}_q$ (combining
    Shannon, Simpson, Gini, Berger--Parker, and McIntosh indices with weights
    $\mu_1, \ldots, \mu_5$).

  \item[Novelty component $N(t,q)$:]
    Penalises tuples whose attribute values are over-represented in
    $\mathcal{H}_q$, favouring attribute–value combinations that have appeared
    sparsely in the session history.
\end{description}

%─────────────────────────────────────────
\subsection{Four Workload Parameters}
\label{sec:theory:params}

The behaviour of each recommender is controlled by four workload parameters
that can be tuned independently in a synthetic benchmark:

\begin{description}
  \item[$\rho$ --- signal fraction:]
    The fraction of tuples in typical large query results that are
    rule-conforming:
    $\rho = \mathbb{E}[|R_q \cap \mathcal{T}^*| / |R_q|]$.
    In our benchmark, $\rho \approx 0.20$--$0.30$ for large queries.

  \item[$\xi$ --- frequency skew:]
    The ratio of the empirical count of common values to the count of rare
    (rule-antecedent) values in $\mathcal{H}_q$:
    $\xi = \text{freq}_q(\text{common}) / \text{freq}_q(\text{rare})$.
    Controlled by the bias probability $p_{\mathrm{common}}$ (set to $0.80$).

  \item[$\sigma^2_{\mathrm{noise}}$ --- noise cluster variance:]
    The mean intra-cluster variance of noise tuples in encoded feature space.
    When $p_{\mathrm{common}}$ is high, the noise distribution concentrates on
    $2$--$3$ values per attribute, making $\sigma^2_{\mathrm{noise}}$ small.

  \item[$Q^*$ --- required memory depth:]
    The minimum session length (number of queries before the evaluation gap)
    needed for MDI to mine association rules with sufficient frequency to
    produce reliable confidence estimates.  Depends on $\text{supp}(r)$ and
    the minimum support threshold $\theta_s$.
\end{description}

%─────────────────────────────────────────
\subsection{Failure Conditions for Baseline Recommenders}
\label{sec:theory:failures}

We now state conditions under which each baseline fails to identify
rule-conforming overlap tuples, giving MDI a ranking advantage.

\begin{proposition}[Frequency failure]
\label{prop:freq}
Let $f_c = \text{freq}_q(\text{common})$ and $f_r = \text{freq}_q(\text{rare})$
be the mean historical counts of common and rare attribute values respectively.
Fix a query result $R_q$ containing noise tuples $R_q^{\rm noise}$ biased toward
common values and rule tuples $R_q^* \subseteq \mathcal{T}^*$.
If $\xi = f_c / f_r > 1$, then for every rule tuple $t^* \in R_q^*$ and typical
noise tuple $t^n \in R_q^{\rm noise}$,
\[
  s_{\mathrm{freq}}(t^n) > s_{\mathrm{freq}}(t^*).
\]
\textit{Proof sketch.}
By construction, each component $\text{freq}_q(A_i = t^n(A_i)) = f_c$ (biased
toward common values) while $\text{freq}_q(A_i = t^*(A_i)) = f_r$ for the
rule-specific attributes.  The inequality follows directly from Eq.~\eqref{eq:freq-score}
and $\xi > 1$.
\end{proposition}

\begin{proposition}[Clustering failure]
\label{prop:clust}
Suppose the noise distribution concentrates on $c$ values per attribute
($c \ll |\text{dom}(A_i)|$) so that $\sigma^2_{\mathrm{noise}} \leq \delta$
for a small $\delta > 0$, and let rule tuples scatter uniformly across
$|\text{dom}(A_i)|$ values.  Then for any $K$-means solution with $K$ clusters,
at least $K - O(\sqrt{\delta})$ centroids coincide with noise cluster means, and
every rule tuple $t^*$ receives an assignment distance
$\|{\phi}(t^*) - \mu_{k(t^*)}\|^2 \geq \Omega(1/c)$, while
$\|{\phi}(t^n) - \mu_{k(t^n)}\|^2 \leq \delta$ for noise tuples.
By Eq.~\eqref{eq:clust-score}, $s_{\mathrm{clust}}(t^n) > s_{\mathrm{clust}}(t^*)$.
\end{proposition}

\begin{proposition}[Similarity failure]
\label{prop:sim}
When $|R_q^{\rm noise}| \gg |R_q^*|$, the result centroid
$\bar{\phi}_q$ is dominated by noise encodings.  Since rule tuples use
attribute values orthogonal to common values in the one-hot representation,
$\phi(t^*) \cdot \bar{\phi}_q \approx \rho\,\|\phi(t^*)\|$, whereas
$\phi(t^n) \cdot \bar{\phi}_q \approx (1-\rho)\,\|\phi(t^n)\|$.  For
$\rho < 1/2$, the cosine score satisfies $s_{\mathrm{sim}}(t^n) > s_{\mathrm{sim}}(t^*)$.
\end{proposition}

\begin{proposition}[Random failure]
\label{prop:rand}
An unbiased random ranking of $n = |R_q|$ tuples, of which $m = |R_q^*|$ are
relevant (ground truth), achieves expected nDCG
\[
  \mathbb{E}[\text{nDCG}_k] \leq \frac{m/n \cdot H_k}{Z_k},
\]
where $H_k = \sum_{i=1}^k 1/\log_2(i+1)$ is the un-normalised ideal DCG.
This equals the ideal nDCG only when $m/n = 1$; for $\rho = m/n \ll 1$
(sparse ground truth), random performance degrades proportionally to $\rho$.
\end{proposition}

\paragraph{Remark on Sampling.}
Sampling interpolates between Random and Frequency: with sufficient history it
converges to the Frequency scorer.  Therefore Propositions~\ref{prop:freq}
and~\ref{prop:rand} bound its performance from above and below respectively,
and it fails for the same structural reason as Frequency once $\xi > 1$.

%─────────────────────────────────────────
\subsection{MDI's Temporal Accumulation Advantage}
\label{sec:theory:mdi}

The key distinction of MDI is that its association component $A(t,q)$
(Eq.~\ref{eq:Acomponent}) \emph{grows} with session history while the baselines
remain stationary or degrade.

\begin{proposition}[Temporal accumulation]
\label{prop:temporal}
Let $r^* = (X^* \Rightarrow Y^*)$ be a ground-truth embedded rule with
confidence $c^* = \text{conf}(r^*)$ and support $s^* = \text{supp}(r^*)$.
Suppose each large query in the session contains a \emph{persistent core} of
$m_c$ rule-conforming tuples drawn i.i.d. from $\mathcal{T}^*$, and MDI mines
$r^*$ from accumulated data at step $q$ with estimated confidence
$\hat{c}_q(r^*) \to c^*$ as $q \to \infty$.  Then the association component
for any tuple $t^* \models X^* \cup Y^*$ satisfies
\[
  A(t^*, q) \xrightarrow{q \to \infty}
  \frac{c^*\,(\lambda_1 + \lambda_3 \cdot c^*/s_Y^* + \lambda_4 \cdot J^*)}{|X^* \cup Y^*|},
\]
where $s_Y^* = \text{supp}(Y^*)$ and $J^*$ is the J-measure of $r^*$.
Moreover, since the rule decay weight $w_{r^*}(q) = e^{-\lambda_r(q - q_{r^*})}$
is refreshed every time $r^*$ is re-mined (i.e., at every large query), the
effective decay does not erode the signal over the session.
\end{proposition}

\begin{corollary}[Gap advantage]
\label{cor:gap}
Because $A(t^*, q)$ accumulates across queries regardless of the prediction gap
$g$, while $s_{\mathrm{freq}}(t)$, $s_{\mathrm{clust}}(t)$, and
$s_{\mathrm{sim}}(t)$ are memoryless (computed from $R_q$ and recent history
only), MDI's advantage over these baselines is \emph{non-decreasing} in session
length $q$ and \emph{independent} of $g$.  This means that at large gaps
($g = 5, 10$), MDI retains its advantage while memoryless baselines degrade
because the current result $R_{q+g}$ drifts from $R_q$ in ways that cannot be
predicted from $R_q$ alone.
\end{corollary}

%─────────────────────────────────────────
\subsection{Benchmark Construction: A Prescriptive Recipe}
\label{sec:theory:recipe}

The preceding propositions jointly characterise when MDI wins, but the
framework also prescribes how to \emph{deliberately construct workloads that
favour any given recommender}.  We state this as a theorem and summarise the
parameter regimes in Table~\ref{tab:recipe}.

\begin{theorem}[Workload-determined advantage]
\label{thm:recipe}
Let $\mathcal{B} = \{\text{MDI, Frequency, Clustering, Similarity, Random,
Sampling}\}$ be the recommender set.  For each $B \in \mathcal{B}$ there
exists a workload characterised by parameters $(\rho, \xi, \sigma^2_{\rm noise},
Q^*)$ such that $\text{nDCG}_k(B, \mathcal{S}) > \text{nDCG}_k(B', \mathcal{S})$
for all $B' \neq B$.  The dominating regime for each recommender is:
\begin{itemize}
  \item \textbf{MDI:} $\rho < 1/2$, $\xi \gg 1$,
        $\sigma^2_{\rm noise} \ll 1$, $Q \geq Q^*(s^*, \theta_s)$.
        Rule-conforming overlap tuples are rare, noise is homogeneous, and
        the session is long enough for rule accumulation.

  \item \textbf{Frequency:} $\rho \approx 1$ and $\xi < 1$ (rare values carry the
        overlap signal, but all overlap tuples have above-average frequency).

  \item \textbf{Clustering:} Overlap tuples form a tight, distinct cluster
        ($\sigma^2_{\rm overlap} \ll \sigma^2_{\rm noise}$) separated from noise
        in encoded space; $K \geq 2$.

  \item \textbf{Similarity:} The future result $R_{q+g}$ is \emph{more similar}
        to the current result centroid than any novel tuple; queries are
        monotonic refinements of the same analytical theme.

  \item \textbf{Random / Sampling:} $\rho \approx 1/2$ and the session is short
        ($Q < Q^*$); association rules are not yet reliable, so memory provides
        no signal and random selection is near-optimal.
\end{itemize}
\end{theorem}

\begin{table}[t!]
  \centering
  \caption{%
    Workload parameter regimes that favour each recommender.
    $\checkmark$ = condition required; $\times$ = not required
    (or reversed).
  }
  \label{tab:recipe}
  \footnotesize
  \setlength{\tabcolsep}{3pt}
  \begin{tabular}{l c c c c}
    \toprule
    \textbf{Rec.}
      & \makecell{$\rho \!\ll\! 1$\\[-2pt]\scriptsize sparse sig.}
      & \makecell{$\xi \!\gg\! 1$\\[-2pt]\scriptsize freq.\ skew}
      & \makecell{$\sigma^2_{\rm n} \!\ll\! 1$\\[-2pt]\scriptsize noise hom.}
      & \makecell{$Q \!\geq\! Q^*$\\[-2pt]\scriptsize deep hist.} \\
    \midrule
    MDI        & \checkmark & \checkmark & \checkmark & \checkmark \\
    Frequency  & $\times$   & $\times$   & $\times$   & \checkmark \\
    Clustering & $\times$   & $\times$   & $\times$\textsuperscript{a}  & $\times$ \\
    Similarity & $\times$   & $\times$   & $\times$\textsuperscript{b}  & $\times$ \\
    Random     & $\times$   & $\times$   & $\times$   & $\times$   \\
    Sampling   & $\times$   & $\times$   & $\times$   & $\times$   \\
    \bottomrule
    \multicolumn{5}{l}{\scriptsize \textsuperscript{a}distinct clusters \quad \textsuperscript{b}centroid-aligned}
  \end{tabular}
\end{table}

\paragraph{Instantiation in the benchmark generator.}
Table~\ref{tab:benchmark-params} shows how the parameters of
\texttt{generate\_benchmark.py} map to the theoretical quantities.  All four
conditions in Theorem~\ref{thm:recipe} for MDI are satisfied by construction:
$\rho \approx 0.22$ for large queries; $\xi = p_{\rm common}/(1 - p_{\rm common}) = 4$
($p_{\rm common} = 0.80$); noise rows concentrate on 2 values per attribute in a
9-attribute schema so $\sigma^2_{\rm noise} \leq 0.11$; and with $Q = 30$ queries
per session (each large query contributing $\sim 150$ tuples) the expected
antecedent support count after 10 large queries is
$10 \times 150 \times \text{supp}(r^*) \approx 90 > 1/\theta_s = 12.5$,
satisfying $Q \geq Q^*$.

\begin{table}[t!]
  \centering
  \caption{%
    Mapping from \texttt{generate\_benchmark.py} parameters to theoretical
    quantities.
  }
  \label{tab:benchmark-params}
  \footnotesize
  \setlength{\tabcolsep}{3pt}
  \begin{tabular}{l l l}
    \toprule
    \textbf{Theory param.} & \textbf{Generator param.} & \textbf{Value} \\
    \midrule
    $\rho$ (signal frac.)     & rule rows / lg.\ query sz.       & $\approx 0.22$ \\
    $\xi$ (freq.\ skew)       & $p_{\rm common} = 0.80$ bias     & $4.0$ \\
    $\sigma^2_{\rm noise}$    & 2 common vals / 9 attrs           & $\leq 0.11$ \\
    $\text{supp}(r^*)$        & 12\% rule rows; conf $= 0.85$    & $\approx 0.10$ \\
    $Q^*$                     & min\_supp $= 0.08$, 3000 rows    & $\approx 3$ lg.\ queries \\
    Core overlap $m_c$        & 20 persistent rule tuples         & $20$ \\
    Sm.\ query overlap        & 24--30 of 40--50 tuples           & $\approx 0.60$--$0.70$ \\
    \bottomrule
  \end{tabular}
\end{table}

\paragraph{Implications for real-world workloads.}
Real EDA sessions over operational databases frequently exhibit high $\xi$ simply
because data skew is ubiquitous, but they rarely satisfy $\rho \ll 1$ unless the
analyst is explicitly drilling into a rare sub-population.  This means MDI's
advantage is most pronounced in \emph{drill-down} scenarios --- exactly the kind
of session studied in the SDSS and SIMBA benchmarks --- and shrinks toward the
level of simpler baselines in broad exploratory sweeps where overlap is spread
uniformly across result rows.  The framework thus not only explains the benchmark
results but also signals where MDI deployment will deliver the greatest practical
benefit.

% Place between the framework description and the Research Challenges section.
% ─────────────────────────────────────────────────────────────────────────────

\section{Theoretical Analysis: Workload--Recommender Alignment}
\label{sec:theory}

The purpose of this section is to give the benchmark results a principled
foundation: we first construct a formal
model of EDA sessions and define precisely what each recommender optimises, then
derive conditions on the \emph{workload} under which each recommender succeeds
or fails.  The analysis yields both an explanation of the benchmark design and a
prescriptive recipe for constructing controlled evaluations targeted at any
recommender.

%─────────────────────────────────────────
\subsection{Formal Session Model}
\label{sec:theory:model}

\paragraph{Schema and tuples.}
Let $\mathcal{A} = \{A_1, \ldots, A_d\}$ be a fixed relational schema.  Each
attribute $A_i$ has a finite domain $\text{dom}(A_i)$.  A \emph{tuple}
$t \in \mathcal{T}$ is a complete assignment $t : \mathcal{A} \to
\bigcup_i \text{dom}(A_i)$ such that $t(A_i) \in \text{dom}(A_i)$.  The
\emph{fact table} $\mathcal{T}$ is a finite multiset of such tuples.

\paragraph{Association rules.}
An \emph{association rule} $r = (X \Rightarrow Y)$ is a pair of attribute–value
sets $X, Y \subset \{A_i = v \mid A_i \in \mathcal{A},\, v \in \text{dom}(A_i)\}$
with $X \cap Y = \emptyset$.  A tuple $t$ \emph{satisfies the antecedent} $X$
if $t(A_i) = v$ for all $(A_i = v) \in X$, and similarly for $Y$.  Standard
interestingness measures are:
\begin{align}
  \text{supp}(r) &= \frac{|\{t \in \mathcal{T} \mid t \models X \cup Y\}|}{|\mathcal{T}|},
  \label{eq:support}\\
  \text{conf}(r) &= \frac{\text{supp}(r)}{\text{supp}(X)},
  \label{eq:confidence}\\
  \text{lift}(r) &= \frac{\text{conf}(r)}{\text{supp}(Y)},
  \label{eq:lift}
\end{align}
where $\text{supp}(X)$ (resp. $\text{supp}(Y)$) denotes the fraction of tuples
satisfying $X$ (resp. $Y$) alone.  The \emph{J-measure}~\cite{smyth1992rule}
is a mutual-information proxy:
\begin{equation}
  J(r) = \text{supp}(r)\log\frac{\text{conf}(r)}{\text{supp}(Y)}
        + (1-\text{supp}(r))\log\frac{1-\text{conf}(r)}{1-\text{supp}(Y)}.
  \label{eq:jmeasure}
\end{equation}

\paragraph{EDA session and prediction task.}
An \emph{EDA session} of length $Q$ is an ordered sequence of query results
$\mathcal{S} = (R_0, R_1, \ldots, R_{Q-1})$ where each $R_q \subseteq \mathcal{T}$
is a finite multiset.  The \emph{history} at step $q$ is
$\mathcal{H}_q = (R_0, \ldots, R_{q-1})$.

At each step $q$ with \emph{prediction gap} $g \geq 1$, a recommender is given
$(R_q, \mathcal{H}_q)$ and must rank the tuples of $R_q$.  The ground-truth
relevance label for a tuple $t \in R_q$ is
\begin{equation}
  \text{rel}(t,q,g) = \mathbf{1}\bigl[t \in R_{q+g}\bigr].
  \label{eq:relevance}
\end{equation}
That is, a tuple is relevant if it \emph{reappears} in the result set $g$ steps
into the future --- a natural formalisation of ``the analyst will want to see
this again soon''.

%─────────────────────────────────────────
\subsection{Evaluation: nDCG as the Quality Criterion}
\label{sec:theory:ndcg}

Let $\pi$ be the ranking of $R_q$ produced by a recommender and let $k$ be the
recommendation budget.  The normalised discounted cumulative gain at rank $k$
is
\begin{equation}
  \text{nDCG}_k(\pi, q, g)
  = \frac{1}{Z_k}
    \sum_{i=1}^{k}
    \frac{\text{rel}(\pi(i), q, g)}{\log_2(i+1)},
  \label{eq:ndcg}
\end{equation}
where $Z_k$ is the ideal DCG (all relevant tuples at the top).  We write
$\text{nDCG}_k(B, \mathcal{S})$ for the mean over all valid $(q, g)$ pairs in
session $\mathcal{S}$ for recommender $B$.

\paragraph{Overlap concentration.}
Define the \emph{rule-conforming overlap fraction} of a query pair $(q, q+g)$
as
\begin{equation}
  \rho(q,g)
  = \frac{|R_q \cap R_{q+g} \cap \mathcal{T}^*|}{|R_q \cap R_{q+g}|},
  \label{eq:rho}
\end{equation}
where $\mathcal{T}^* = \{t \in \mathcal{T} \mid \exists r : t \models X_r
\cup Y_r\}$ is the set of rule-conforming tuples.  When $\rho(q,g) \approx 1$,
predicting overlap is essentially the same as identifying rule tuples inside
$R_q$.

%─────────────────────────────────────────
\subsection{Recommender Scoring Functions}
\label{sec:theory:scorers}

We model each recommender as a \emph{scoring function}
$s_B : R_q \times \mathcal{H}_q \to \mathbb{R}$ that induces a ranking $\pi_B$
over $R_q$.

\paragraph{Random.}
$s_{\mathrm{rand}}(t) \sim \mathrm{Uniform}(0,1)$, independent of everything.

\paragraph{Frequency.}
Let $\text{freq}_q(A_i = v) = |\{(s, u) \mid s < q,\, u \in R_s,\, u(A_i) = v\}|$
be the historical count of value $v$ in attribute $A_i$ across past results.
\begin{equation}
  s_{\mathrm{freq}}(t) = \frac{1}{d}\sum_{i=1}^d \text{freq}_q\bigl(A_i = t(A_i)\bigr).
  \label{eq:freq-score}
\end{equation}

\paragraph{Clustering.}
One-hot encode each tuple in $R_q$ to a vector $\phi(t) \in \{0,1\}^D$.  Run
$K$-means on $\{\phi(t) \mid t \in R_q\}$ to obtain centroids
$\{\mu_1, \ldots, \mu_K\}$.  Score each tuple by proximity to its assigned
centroid's \emph{neighbourhood density}:
\begin{equation}
  s_{\mathrm{clust}}(t) = -\|\phi(t) - \mu_{k(t)}\|^2,
  \label{eq:clust-score}
\end{equation}
where $k(t) = \arg\min_k \|\phi(t) - \mu_k\|$.  Tuples near a centroid score
high; outliers score low.

\paragraph{Similarity.}
Let $\bar{\phi}_q = \frac{1}{|R_q|}\sum_{t \in R_q}\phi(t)$ be the mean
encoding of the current result.  Rank by cosine similarity to this centroid:
\begin{equation}
  s_{\mathrm{sim}}(t) = \frac{\phi(t) \cdot \bar{\phi}_q}{\|\phi(t)\|\,\|\bar{\phi}_q\|}.
  \label{eq:sim-score}
\end{equation}

\paragraph{Sampling.}
A stratified random sample of $k$ tuples from $R_q \cup \mathcal{H}_{q}^{\rm flat}$
(the union of current and past results), with each tuple weighted by its
historical frequency.  This interpolates between Random and Frequency depending
on the history length.

\paragraph{MDI.}
The MDI score accumulates three components across the history:
\begin{equation}
  s_{\mathrm{MDI}}(t, q)
  = \alpha \cdot A(t, q) + \beta \cdot D(t, q) + \gamma \cdot N(t, q),
  \label{eq:mdi-score}
\end{equation}
where $\alpha + \beta + \gamma = 1$.

\begin{description}
  \item[Association component $A(t,q)$:]
    Let $\mathcal{R}_q$ be association rules mined from the history up to step
    $q$.  For rule $r \in \mathcal{R}_q$, let
    $w_r(q) = e^{-\lambda_r (q - q_r)}$ be an exponential decay weight
    ($\lambda_r$ is the decay rate, $q_r$ the query at which $r$ was mined).
    The composite rule score is
    \begin{equation}
      I_{\mathrm{assoc}}(r)
      = \lambda_1\,\mathrm{conf}(r) + \lambda_2\,\mathrm{supp}(r)
      + \lambda_3\,\mathrm{lift}(r) + \lambda_4\,J(r),
      \label{eq:Iassoc}
    \end{equation}
    and the association component for tuple $t$ is
    \begin{equation}
      A(t,q)
      = \sum_{r \in \mathcal{R}_q}
        w_r(q)\,\frac{I_{\mathrm{assoc}}(r)}{|X_r \cup Y_r|}\,
        \mathbf{1}[t \models X_r \cup Y_r].
      \label{eq:Acomponent}
    \end{equation}

  \item[Diversity component $D(t,q)$:]
    Measures how much adding $t$ to a candidate summary would increase the
    entropy of attribute-value distributions in $\mathcal{H}_q$ (combining
    Shannon, Simpson, Gini, Berger--Parker, and McIntosh indices with weights
    $\mu_1, \ldots, \mu_5$).

  \item[Novelty component $N(t,q)$:]
    Penalises tuples whose attribute values are over-represented in
    $\mathcal{H}_q$, favouring attribute–value combinations that have appeared
    sparsely in the session history.
\end{description}

%─────────────────────────────────────────
\subsection{Four Workload Parameters}
\label{sec:theory:params}

The behaviour of each recommender is controlled by four workload parameters
that can be tuned independently in a synthetic benchmark:

\begin{description}
  \item[$\rho$ --- signal fraction:]
    The fraction of tuples in typical large query results that are
    rule-conforming:
    $\rho = \mathbb{E}[|R_q \cap \mathcal{T}^*| / |R_q|]$.
    In our benchmark, $\rho \approx 0.20$--$0.30$ for large queries.

  \item[$\xi$ --- frequency skew:]
    The ratio of the empirical count of common values to the count of rare
    (rule-antecedent) values in $\mathcal{H}_q$:
    $\xi = \text{freq}_q(\text{common}) / \text{freq}_q(\text{rare})$.
    Controlled by the bias probability $p_{\mathrm{common}}$ (set to $0.80$).

  \item[$\sigma^2_{\mathrm{noise}}$ --- noise cluster variance:]
    The mean intra-cluster variance of noise tuples in encoded feature space.
    When $p_{\mathrm{common}}$ is high, the noise distribution concentrates on
    $2$--$3$ values per attribute, making $\sigma^2_{\mathrm{noise}}$ small.

  \item[$Q^*$ --- required memory depth:]
    The minimum session length (number of queries before the evaluation gap)
    needed for MDI to mine association rules with sufficient frequency to
    produce reliable confidence estimates.  Depends on $\text{supp}(r)$ and
    the minimum support threshold $\theta_s$.
\end{description}

%─────────────────────────────────────────
\subsection{Failure Conditions for Baseline Recommenders}
\label{sec:theory:failures}

We now state conditions under which each baseline fails to identify
rule-conforming overlap tuples, giving MDI a ranking advantage.

\begin{proposition}[Frequency failure]
\label{prop:freq}
Let $f_c = \text{freq}_q(\text{common})$ and $f_r = \text{freq}_q(\text{rare})$
be the mean historical counts of common and rare attribute values respectively.
Fix a query result $R_q$ containing noise tuples $R_q^{\rm noise}$ biased toward
common values and rule tuples $R_q^* \subseteq \mathcal{T}^*$.
If $\xi = f_c / f_r > 1$, then for every rule tuple $t^* \in R_q^*$ and typical
noise tuple $t^n \in R_q^{\rm noise}$,
\[
  s_{\mathrm{freq}}(t^n) > s_{\mathrm{freq}}(t^*).
\]
\textit{Proof sketch.}
By construction, each component $\text{freq}_q(A_i = t^n(A_i)) = f_c$ (biased
toward common values) while $\text{freq}_q(A_i = t^*(A_i)) = f_r$ for the
rule-specific attributes.  The inequality follows directly from Eq.~\eqref{eq:freq-score}
and $\xi > 1$.
\end{proposition}

\begin{proposition}[Clustering failure]
\label{prop:clust}
Suppose the noise distribution concentrates on $c$ values per attribute
($c \ll |\text{dom}(A_i)|$) so that $\sigma^2_{\mathrm{noise}} \leq \delta$
for a small $\delta > 0$, and let rule tuples scatter uniformly across
$|\text{dom}(A_i)|$ values.  Then for any $K$-means solution with $K$ clusters,
at least $K - O(\sqrt{\delta})$ centroids coincide with noise cluster means, and
every rule tuple $t^*$ receives an assignment distance
$\|{\phi}(t^*) - \mu_{k(t^*)}\|^2 \geq \Omega(1/c)$, while
$\|{\phi}(t^n) - \mu_{k(t^n)}\|^2 \leq \delta$ for noise tuples.
By Eq.~\eqref{eq:clust-score}, $s_{\mathrm{clust}}(t^n) > s_{\mathrm{clust}}(t^*)$.
\end{proposition}

\begin{proposition}[Similarity failure]
\label{prop:sim}
When $|R_q^{\rm noise}| \gg |R_q^*|$, the result centroid
$\bar{\phi}_q$ is dominated by noise encodings.  Since rule tuples use
attribute values orthogonal to common values in the one-hot representation,
$\phi(t^*) \cdot \bar{\phi}_q \approx \rho\,\|\phi(t^*)\|$, whereas
$\phi(t^n) \cdot \bar{\phi}_q \approx (1-\rho)\,\|\phi(t^n)\|$.  For
$\rho < 1/2$, the cosine score satisfies $s_{\mathrm{sim}}(t^n) > s_{\mathrm{sim}}(t^*)$.
\end{proposition}

\begin{proposition}[Random failure]
\label{prop:rand}
An unbiased random ranking of $n = |R_q|$ tuples, of which $m = |R_q^*|$ are
relevant (ground truth), achieves expected nDCG
\[
  \mathbb{E}[\text{nDCG}_k] \leq \frac{m/n \cdot H_k}{Z_k},
\]
where $H_k = \sum_{i=1}^k 1/\log_2(i+1)$ is the un-normalised ideal DCG.
This equals the ideal nDCG only when $m/n = 1$; for $\rho = m/n \ll 1$
(sparse ground truth), random performance degrades proportionally to $\rho$.
\end{proposition}

\paragraph{Remark on Sampling.}
Sampling interpolates between Random and Frequency: with sufficient history it
converges to the Frequency scorer.  Therefore Propositions~\ref{prop:freq}
and~\ref{prop:rand} bound its performance from above and below respectively,
and it fails for the same structural reason as Frequency once $\xi > 1$.

%─────────────────────────────────────────
\subsection{MDI's Temporal Accumulation Advantage}
\label{sec:theory:mdi}

The key distinction of MDI is that its association component $A(t,q)$
(Eq.~\ref{eq:Acomponent}) \emph{grows} with session history while the baselines
remain stationary or degrade.

\begin{proposition}[Temporal accumulation]
\label{prop:temporal}
Let $r^* = (X^* \Rightarrow Y^*)$ be a ground-truth embedded rule with
confidence $c^* = \text{conf}(r^*)$ and support $s^* = \text{supp}(r^*)$.
Suppose each large query in the session contains a \emph{persistent core} of
$m_c$ rule-conforming tuples drawn i.i.d. from $\mathcal{T}^*$, and MDI mines
$r^*$ from accumulated data at step $q$ with estimated confidence
$\hat{c}_q(r^*) \to c^*$ as $q \to \infty$.  Then the association component
for any tuple $t^* \models X^* \cup Y^*$ satisfies
\[
  A(t^*, q) \xrightarrow{q \to \infty}
  \frac{c^*\,(\lambda_1 + \lambda_3 \cdot c^*/s_Y^* + \lambda_4 \cdot J^*)}{|X^* \cup Y^*|},
\]
where $s_Y^* = \text{supp}(Y^*)$ and $J^*$ is the J-measure of $r^*$.
Moreover, since the rule decay weight $w_{r^*}(q) = e^{-\lambda_r(q - q_{r^*})}$
is refreshed every time $r^*$ is re-mined (i.e., at every large query), the
effective decay does not erode the signal over the session.
\end{proposition}

\begin{corollary}[Gap advantage]
\label{cor:gap}
Because $A(t^*, q)$ accumulates across queries regardless of the prediction gap
$g$, while $s_{\mathrm{freq}}(t)$, $s_{\mathrm{clust}}(t)$, and
$s_{\mathrm{sim}}(t)$ are memoryless (computed from $R_q$ and recent history
only), MDI's advantage over these baselines is \emph{non-decreasing} in session
length $q$ and \emph{independent} of $g$.  This means that at large gaps
($g = 5, 10$), MDI retains its advantage while memoryless baselines degrade
because the current result $R_{q+g}$ drifts from $R_q$ in ways that cannot be
predicted from $R_q$ alone.
\end{corollary}

%─────────────────────────────────────────
\subsection{Benchmark Construction: A Prescriptive Recipe}
\label{sec:theory:recipe}

The preceding propositions jointly characterise when MDI wins, but the
framework also prescribes how to \emph{deliberately construct workloads that
favour any given recommender}.  We state this as a theorem and summarise the
parameter regimes in Table~\ref{tab:recipe}.

\begin{theorem}[Workload-determined advantage]
\label{thm:recipe}
Let $\mathcal{B} = \{\text{MDI, Frequency, Clustering, Similarity, Random,
Sampling}\}$ be the recommender set.  For each $B \in \mathcal{B}$ there
exists a workload characterised by parameters $(\rho, \xi, \sigma^2_{\rm noise},
Q^*)$ such that $\text{nDCG}_k(B, \mathcal{S}) > \text{nDCG}_k(B', \mathcal{S})$
for all $B' \neq B$.  The dominating regime for each recommender is:
\begin{itemize}
  \item \textbf{MDI:} $\rho < 1/2$, $\xi \gg 1$,
        $\sigma^2_{\rm noise} \ll 1$, $Q \geq Q^*(s^*, \theta_s)$.
        Rule-conforming overlap tuples are rare, noise is homogeneous, and
        the session is long enough for rule accumulation.

  \item \textbf{Frequency:} $\rho \approx 1$ and $\xi < 1$ (rare values carry the
        overlap signal, but all overlap tuples have above-average frequency).

  \item \textbf{Clustering:} Overlap tuples form a tight, distinct cluster
        ($\sigma^2_{\rm overlap} \ll \sigma^2_{\rm noise}$) separated from noise
        in encoded space; $K \geq 2$.

  \item \textbf{Similarity:} The future result $R_{q+g}$ is \emph{more similar}
        to the current result centroid than any novel tuple; queries are
        monotonic refinements of the same analytical theme.

  \item \textbf{Random / Sampling:} $\rho \approx 1/2$ and the session is short
        ($Q < Q^*$); association rules are not yet reliable, so memory provides
        no signal and random selection is near-optimal.
\end{itemize}
\end{theorem}

\begin{table}[t!]
  \centering
  \caption{%
    Workload parameter regimes that favour each recommender.
    $\checkmark$ = condition required; $\times$ = not required
    (or reversed).
  }
  \label{tab:recipe}
  \footnotesize
  \setlength{\tabcolsep}{3pt}
  \begin{tabular}{l c c c c}
    \toprule
    \textbf{Rec.}
      & \makecell{$\rho \!\ll\! 1$\\[-2pt]\scriptsize sparse sig.}
      & \makecell{$\xi \!\gg\! 1$\\[-2pt]\scriptsize freq.\ skew}
      & \makecell{$\sigma^2_{\rm n} \!\ll\! 1$\\[-2pt]\scriptsize noise hom.}
      & \makecell{$Q \!\geq\! Q^*$\\[-2pt]\scriptsize deep hist.} \\
    \midrule
    MDI        & \checkmark & \checkmark & \checkmark & \checkmark \\
    Frequency  & $\times$   & $\times$   & $\times$   & \checkmark \\
    Clustering & $\times$   & $\times$   & $\times$\textsuperscript{a}  & $\times$ \\
    Similarity & $\times$   & $\times$   & $\times$\textsuperscript{b}  & $\times$ \\
    Random     & $\times$   & $\times$   & $\times$   & $\times$   \\
    Sampling   & $\times$   & $\times$   & $\times$   & $\times$   \\
    \bottomrule
    \multicolumn{5}{l}{\scriptsize \textsuperscript{a}distinct clusters \quad \textsuperscript{b}centroid-aligned}
  \end{tabular}
\end{table}

\paragraph{Instantiation in the benchmark generator.}
Table~\ref{tab:benchmark-params} shows how the parameters of
\texttt{generate\_benchmark.py} map to the theoretical quantities.  All four
conditions in Theorem~\ref{thm:recipe} for MDI are satisfied by construction:
$\rho \approx 0.22$ for large queries; $\xi = p_{\rm common}/(1 - p_{\rm common}) = 4$
($p_{\rm common} = 0.80$); noise rows concentrate on 2 values per attribute in a
9-attribute schema so $\sigma^2_{\rm noise} \leq 0.11$; and with $Q = 30$ queries
per session (each large query contributing $\sim 150$ tuples) the expected
antecedent support count after 10 large queries is
$10 \times 150 \times \text{supp}(r^*) \approx 90 > 1/\theta_s = 12.5$,
satisfying $Q \geq Q^*$.

\begin{table}[t!]
  \centering
  \caption{%
    Mapping from \texttt{generate\_benchmark.py} parameters to theoretical
    quantities.
  }
  \label{tab:benchmark-params}
  \footnotesize
  \setlength{\tabcolsep}{3pt}
  \begin{tabular}{l l l}
    \toprule
    \textbf{Theory param.} & \textbf{Generator param.} & \textbf{Value} \\
    \midrule
    $\rho$ (signal frac.)     & rule rows / lg.\ query sz.       & $\approx 0.22$ \\
    $\xi$ (freq.\ skew)       & $p_{\rm common} = 0.80$ bias     & $4.0$ \\
    $\sigma^2_{\rm noise}$    & 2 common vals / 9 attrs           & $\leq 0.11$ \\
    $\text{supp}(r^*)$        & 12\% rule rows; conf $= 0.85$    & $\approx 0.10$ \\
    $Q^*$                     & min\_supp $= 0.08$, 3000 rows    & $\approx 3$ lg.\ queries \\
    Core overlap $m_c$        & 20 persistent rule tuples         & $20$ \\
    Sm.\ query overlap        & 24--30 of 40--50 tuples           & $\approx 0.60$--$0.70$ \\
    \bottomrule
  \end{tabular}
\end{table}

\paragraph{Implications for real-world workloads.}
Real EDA sessions over operational databases frequently exhibit high $\xi$ simply
because data skew is ubiquitous, but they rarely satisfy $\rho \ll 1$ unless the
analyst is explicitly drilling into a rare sub-population.  This means MDI's
advantage is most pronounced in \emph{drill-down} scenarios --- exactly the kind
of session studied in the SDSS and SIMBA benchmarks --- and shrinks toward the
level of simpler baselines in broad exploratory sweeps where overlap is spread
uniformly across result rows.  The framework thus not only explains the benchmark
results but also signals where MDI deployment will deliver the greatest practical
benefit.

% Place between the framework description and the Research Challenges section.
% ─────────────────────────────────────────────────────────────────────────────

\section{Theoretical Analysis: Workload--Recommender Alignment}
\label{sec:theory}

The purpose of this section is to give the benchmark results a principled
foundation: we first construct a formal
model of EDA sessions and define precisely what each recommender optimises, then
derive conditions on the \emph{workload} under which each recommender succeeds
or fails.  The analysis yields both an explanation of the benchmark design and a
prescriptive recipe for constructing controlled evaluations targeted at any
recommender.

%─────────────────────────────────────────
\subsection{Formal Session Model}
\label{sec:theory:model}

\paragraph{Schema and tuples.}
Let $\mathcal{A} = \{A_1, \ldots, A_d\}$ be a fixed relational schema.  Each
attribute $A_i$ has a finite domain $\text{dom}(A_i)$.  A \emph{tuple}
$t \in \mathcal{T}$ is a complete assignment $t : \mathcal{A} \to
\bigcup_i \text{dom}(A_i)$ such that $t(A_i) \in \text{dom}(A_i)$.  The
\emph{fact table} $\mathcal{T}$ is a finite multiset of such tuples.

\paragraph{Association rules.}
An \emph{association rule} $r = (X \Rightarrow Y)$ is a pair of attribute–value
sets $X, Y \subset \{A_i = v \mid A_i \in \mathcal{A},\, v \in \text{dom}(A_i)\}$
with $X \cap Y = \emptyset$.  A tuple $t$ \emph{satisfies the antecedent} $X$
if $t(A_i) = v$ for all $(A_i = v) \in X$, and similarly for $Y$.  Standard
interestingness measures are:
\begin{align}
  \text{supp}(r) &= \frac{|\{t \in \mathcal{T} \mid t \models X \cup Y\}|}{|\mathcal{T}|},
  \label{eq:support}\\
  \text{conf}(r) &= \frac{\text{supp}(r)}{\text{supp}(X)},
  \label{eq:confidence}\\
  \text{lift}(r) &= \frac{\text{conf}(r)}{\text{supp}(Y)},
  \label{eq:lift}
\end{align}
where $\text{supp}(X)$ (resp. $\text{supp}(Y)$) denotes the fraction of tuples
satisfying $X$ (resp. $Y$) alone.  The \emph{J-measure}~\cite{smyth1992rule}
is a mutual-information proxy:
\begin{equation}
  J(r) = \text{supp}(r)\log\frac{\text{conf}(r)}{\text{supp}(Y)}
        + (1-\text{supp}(r))\log\frac{1-\text{conf}(r)}{1-\text{supp}(Y)}.
  \label{eq:jmeasure}
\end{equation}

\paragraph{EDA session and prediction task.}
An \emph{EDA session} of length $Q$ is an ordered sequence of query results
$\mathcal{S} = (R_0, R_1, \ldots, R_{Q-1})$ where each $R_q \subseteq \mathcal{T}$
is a finite multiset.  The \emph{history} at step $q$ is
$\mathcal{H}_q = (R_0, \ldots, R_{q-1})$.

At each step $q$ with \emph{prediction gap} $g \geq 1$, a recommender is given
$(R_q, \mathcal{H}_q)$ and must rank the tuples of $R_q$.  The ground-truth
relevance label for a tuple $t \in R_q$ is
\begin{equation}
  \text{rel}(t,q,g) = \mathbf{1}\bigl[t \in R_{q+g}\bigr].
  \label{eq:relevance}
\end{equation}
That is, a tuple is relevant if it \emph{reappears} in the result set $g$ steps
into the future --- a natural formalisation of ``the analyst will want to see
this again soon''.

%─────────────────────────────────────────
\subsection{Evaluation: nDCG as the Quality Criterion}
\label{sec:theory:ndcg}

Let $\pi$ be the ranking of $R_q$ produced by a recommender and let $k$ be the
recommendation budget.  The normalised discounted cumulative gain at rank $k$
is
\begin{equation}
  \text{nDCG}_k(\pi, q, g)
  = \frac{1}{Z_k}
    \sum_{i=1}^{k}
    \frac{\text{rel}(\pi(i), q, g)}{\log_2(i+1)},
  \label{eq:ndcg}
\end{equation}
where $Z_k$ is the ideal DCG (all relevant tuples at the top).  We write
$\text{nDCG}_k(B, \mathcal{S})$ for the mean over all valid $(q, g)$ pairs in
session $\mathcal{S}$ for recommender $B$.

\paragraph{Overlap concentration.}
Define the \emph{rule-conforming overlap fraction} of a query pair $(q, q+g)$
as
\begin{equation}
  \rho(q,g)
  = \frac{|R_q \cap R_{q+g} \cap \mathcal{T}^*|}{|R_q \cap R_{q+g}|},
  \label{eq:rho}
\end{equation}
where $\mathcal{T}^* = \{t \in \mathcal{T} \mid \exists r : t \models X_r
\cup Y_r\}$ is the set of rule-conforming tuples.  When $\rho(q,g) \approx 1$,
predicting overlap is essentially the same as identifying rule tuples inside
$R_q$.

%─────────────────────────────────────────
\subsection{Recommender Scoring Functions}
\label{sec:theory:scorers}

We model each recommender as a \emph{scoring function}
$s_B : R_q \times \mathcal{H}_q \to \mathbb{R}$ that induces a ranking $\pi_B$
over $R_q$.

\paragraph{Random.}
$s_{\mathrm{rand}}(t) \sim \mathrm{Uniform}(0,1)$, independent of everything.

\paragraph{Frequency.}
Let $\text{freq}_q(A_i = v) = |\{(s, u) \mid s < q,\, u \in R_s,\, u(A_i) = v\}|$
be the historical count of value $v$ in attribute $A_i$ across past results.
\begin{equation}
  s_{\mathrm{freq}}(t) = \frac{1}{d}\sum_{i=1}^d \text{freq}_q\bigl(A_i = t(A_i)\bigr).
  \label{eq:freq-score}
\end{equation}

\paragraph{Clustering.}
One-hot encode each tuple in $R_q$ to a vector $\phi(t) \in \{0,1\}^D$.  Run
$K$-means on $\{\phi(t) \mid t \in R_q\}$ to obtain centroids
$\{\mu_1, \ldots, \mu_K\}$.  Score each tuple by proximity to its assigned
centroid's \emph{neighbourhood density}:
\begin{equation}
  s_{\mathrm{clust}}(t) = -\|\phi(t) - \mu_{k(t)}\|^2,
  \label{eq:clust-score}
\end{equation}
where $k(t) = \arg\min_k \|\phi(t) - \mu_k\|$.  Tuples near a centroid score
high; outliers score low.

\paragraph{Similarity.}
Let $\bar{\phi}_q = \frac{1}{|R_q|}\sum_{t \in R_q}\phi(t)$ be the mean
encoding of the current result.  Rank by cosine similarity to this centroid:
\begin{equation}
  s_{\mathrm{sim}}(t) = \frac{\phi(t) \cdot \bar{\phi}_q}{\|\phi(t)\|\,\|\bar{\phi}_q\|}.
  \label{eq:sim-score}
\end{equation}

\paragraph{Sampling.}
A stratified random sample of $k$ tuples from $R_q \cup \mathcal{H}_{q}^{\rm flat}$
(the union of current and past results), with each tuple weighted by its
historical frequency.  This interpolates between Random and Frequency depending
on the history length.

\paragraph{MDI.}
The MDI score accumulates three components across the history:
\begin{equation}
  s_{\mathrm{MDI}}(t, q)
  = \alpha \cdot A(t, q) + \beta \cdot D(t, q) + \gamma \cdot N(t, q),
  \label{eq:mdi-score}
\end{equation}
where $\alpha + \beta + \gamma = 1$.

\begin{description}
  \item[Association component $A(t,q)$:]
    Let $\mathcal{R}_q$ be association rules mined from the history up to step
    $q$.  For rule $r \in \mathcal{R}_q$, let
    $w_r(q) = e^{-\lambda_r (q - q_r)}$ be an exponential decay weight
    ($\lambda_r$ is the decay rate, $q_r$ the query at which $r$ was mined).
    The composite rule score is
    \begin{equation}
      I_{\mathrm{assoc}}(r)
      = \lambda_1\,\mathrm{conf}(r) + \lambda_2\,\mathrm{supp}(r)
      + \lambda_3\,\mathrm{lift}(r) + \lambda_4\,J(r),
      \label{eq:Iassoc}
    \end{equation}
    and the association component for tuple $t$ is
    \begin{equation}
      A(t,q)
      = \sum_{r \in \mathcal{R}_q}
        w_r(q)\,\frac{I_{\mathrm{assoc}}(r)}{|X_r \cup Y_r|}\,
        \mathbf{1}[t \models X_r \cup Y_r].
      \label{eq:Acomponent}
    \end{equation}

  \item[Diversity component $D(t,q)$:]
    Measures how much adding $t$ to a candidate summary would increase the
    entropy of attribute-value distributions in $\mathcal{H}_q$ (combining
    Shannon, Simpson, Gini, Berger--Parker, and McIntosh indices with weights
    $\mu_1, \ldots, \mu_5$).

  \item[Novelty component $N(t,q)$:]
    Penalises tuples whose attribute values are over-represented in
    $\mathcal{H}_q$, favouring attribute–value combinations that have appeared
    sparsely in the session history.
\end{description}

%─────────────────────────────────────────
\subsection{Four Workload Parameters}
\label{sec:theory:params}

The behaviour of each recommender is controlled by four workload parameters
that can be tuned independently in a synthetic benchmark:

\begin{description}
  \item[$\rho$ --- signal fraction:]
    The fraction of tuples in typical large query results that are
    rule-conforming:
    $\rho = \mathbb{E}[|R_q \cap \mathcal{T}^*| / |R_q|]$.
    In our benchmark, $\rho \approx 0.20$--$0.30$ for large queries.

  \item[$\xi$ --- frequency skew:]
    The ratio of the empirical count of common values to the count of rare
    (rule-antecedent) values in $\mathcal{H}_q$:
    $\xi = \text{freq}_q(\text{common}) / \text{freq}_q(\text{rare})$.
    Controlled by the bias probability $p_{\mathrm{common}}$ (set to $0.80$).

  \item[$\sigma^2_{\mathrm{noise}}$ --- noise cluster variance:]
    The mean intra-cluster variance of noise tuples in encoded feature space.
    When $p_{\mathrm{common}}$ is high, the noise distribution concentrates on
    $2$--$3$ values per attribute, making $\sigma^2_{\mathrm{noise}}$ small.

  \item[$Q^*$ --- required memory depth:]
    The minimum session length (number of queries before the evaluation gap)
    needed for MDI to mine association rules with sufficient frequency to
    produce reliable confidence estimates.  Depends on $\text{supp}(r)$ and
    the minimum support threshold $\theta_s$.
\end{description}

%─────────────────────────────────────────
\subsection{Failure Conditions for Baseline Recommenders}
\label{sec:theory:failures}

We now state conditions under which each baseline fails to identify
rule-conforming overlap tuples, giving MDI a ranking advantage.

\begin{proposition}[Frequency failure]
\label{prop:freq}
Let $f_c = \text{freq}_q(\text{common})$ and $f_r = \text{freq}_q(\text{rare})$
be the mean historical counts of common and rare attribute values respectively.
Fix a query result $R_q$ containing noise tuples $R_q^{\rm noise}$ biased toward
common values and rule tuples $R_q^* \subseteq \mathcal{T}^*$.
If $\xi = f_c / f_r > 1$, then for every rule tuple $t^* \in R_q^*$ and typical
noise tuple $t^n \in R_q^{\rm noise}$,
\[
  s_{\mathrm{freq}}(t^n) > s_{\mathrm{freq}}(t^*).
\]
\textit{Proof sketch.}
By construction, each component $\text{freq}_q(A_i = t^n(A_i)) = f_c$ (biased
toward common values) while $\text{freq}_q(A_i = t^*(A_i)) = f_r$ for the
rule-specific attributes.  The inequality follows directly from Eq.~\eqref{eq:freq-score}
and $\xi > 1$.
\end{proposition}

\begin{proposition}[Clustering failure]
\label{prop:clust}
Suppose the noise distribution concentrates on $c$ values per attribute
($c \ll |\text{dom}(A_i)|$) so that $\sigma^2_{\mathrm{noise}} \leq \delta$
for a small $\delta > 0$, and let rule tuples scatter uniformly across
$|\text{dom}(A_i)|$ values.  Then for any $K$-means solution with $K$ clusters,
at least $K - O(\sqrt{\delta})$ centroids coincide with noise cluster means, and
every rule tuple $t^*$ receives an assignment distance
$\|{\phi}(t^*) - \mu_{k(t^*)}\|^2 \geq \Omega(1/c)$, while
$\|{\phi}(t^n) - \mu_{k(t^n)}\|^2 \leq \delta$ for noise tuples.
By Eq.~\eqref{eq:clust-score}, $s_{\mathrm{clust}}(t^n) > s_{\mathrm{clust}}(t^*)$.
\end{proposition}

\begin{proposition}[Similarity failure]
\label{prop:sim}
When $|R_q^{\rm noise}| \gg |R_q^*|$, the result centroid
$\bar{\phi}_q$ is dominated by noise encodings.  Since rule tuples use
attribute values orthogonal to common values in the one-hot representation,
$\phi(t^*) \cdot \bar{\phi}_q \approx \rho\,\|\phi(t^*)\|$, whereas
$\phi(t^n) \cdot \bar{\phi}_q \approx (1-\rho)\,\|\phi(t^n)\|$.  For
$\rho < 1/2$, the cosine score satisfies $s_{\mathrm{sim}}(t^n) > s_{\mathrm{sim}}(t^*)$.
\end{proposition}

\begin{proposition}[Random failure]
\label{prop:rand}
An unbiased random ranking of $n = |R_q|$ tuples, of which $m = |R_q^*|$ are
relevant (ground truth), achieves expected nDCG
\[
  \mathbb{E}[\text{nDCG}_k] \leq \frac{m/n \cdot H_k}{Z_k},
\]
where $H_k = \sum_{i=1}^k 1/\log_2(i+1)$ is the un-normalised ideal DCG.
This equals the ideal nDCG only when $m/n = 1$; for $\rho = m/n \ll 1$
(sparse ground truth), random performance degrades proportionally to $\rho$.
\end{proposition}

\paragraph{Remark on Sampling.}
Sampling interpolates between Random and Frequency: with sufficient history it
converges to the Frequency scorer.  Therefore Propositions~\ref{prop:freq}
and~\ref{prop:rand} bound its performance from above and below respectively,
and it fails for the same structural reason as Frequency once $\xi > 1$.

%─────────────────────────────────────────
\subsection{MDI's Temporal Accumulation Advantage}
\label{sec:theory:mdi}

The key distinction of MDI is that its association component $A(t,q)$
(Eq.~\ref{eq:Acomponent}) \emph{grows} with session history while the baselines
remain stationary or degrade.

\begin{proposition}[Temporal accumulation]
\label{prop:temporal}
Let $r^* = (X^* \Rightarrow Y^*)$ be a ground-truth embedded rule with
confidence $c^* = \text{conf}(r^*)$ and support $s^* = \text{supp}(r^*)$.
Suppose each large query in the session contains a \emph{persistent core} of
$m_c$ rule-conforming tuples drawn i.i.d. from $\mathcal{T}^*$, and MDI mines
$r^*$ from accumulated data at step $q$ with estimated confidence
$\hat{c}_q(r^*) \to c^*$ as $q \to \infty$.  Then the association component
for any tuple $t^* \models X^* \cup Y^*$ satisfies
\[
  A(t^*, q) \xrightarrow{q \to \infty}
  \frac{c^*\,(\lambda_1 + \lambda_3 \cdot c^*/s_Y^* + \lambda_4 \cdot J^*)}{|X^* \cup Y^*|},
\]
where $s_Y^* = \text{supp}(Y^*)$ and $J^*$ is the J-measure of $r^*$.
Moreover, since the rule decay weight $w_{r^*}(q) = e^{-\lambda_r(q - q_{r^*})}$
is refreshed every time $r^*$ is re-mined (i.e., at every large query), the
effective decay does not erode the signal over the session.
\end{proposition}

\begin{corollary}[Gap advantage]
\label{cor:gap}
Because $A(t^*, q)$ accumulates across queries regardless of the prediction gap
$g$, while $s_{\mathrm{freq}}(t)$, $s_{\mathrm{clust}}(t)$, and
$s_{\mathrm{sim}}(t)$ are memoryless (computed from $R_q$ and recent history
only), MDI's advantage over these baselines is \emph{non-decreasing} in session
length $q$ and \emph{independent} of $g$.  This means that at large gaps
($g = 5, 10$), MDI retains its advantage while memoryless baselines degrade
because the current result $R_{q+g}$ drifts from $R_q$ in ways that cannot be
predicted from $R_q$ alone.
\end{corollary}

%─────────────────────────────────────────
\subsection{Benchmark Construction: A Prescriptive Recipe}
\label{sec:theory:recipe}

The preceding propositions jointly characterise when MDI wins, but the
framework also prescribes how to \emph{deliberately construct workloads that
favour any given recommender}.  We state this as a theorem and summarise the
parameter regimes in Table~\ref{tab:recipe}.

\begin{theorem}[Workload-determined advantage]
\label{thm:recipe}
Let $\mathcal{B} = \{\text{MDI, Frequency, Clustering, Similarity, Random,
Sampling}\}$ be the recommender set.  For each $B \in \mathcal{B}$ there
exists a workload characterised by parameters $(\rho, \xi, \sigma^2_{\rm noise},
Q^*)$ such that $\text{nDCG}_k(B, \mathcal{S}) > \text{nDCG}_k(B', \mathcal{S})$
for all $B' \neq B$.  The dominating regime for each recommender is:
\begin{itemize}
  \item \textbf{MDI:} $\rho < 1/2$, $\xi \gg 1$,
        $\sigma^2_{\rm noise} \ll 1$, $Q \geq Q^*(s^*, \theta_s)$.
        Rule-conforming overlap tuples are rare, noise is homogeneous, and
        the session is long enough for rule accumulation.

  \item \textbf{Frequency:} $\rho \approx 1$ and $\xi < 1$ (rare values carry the
        overlap signal, but all overlap tuples have above-average frequency).

  \item \textbf{Clustering:} Overlap tuples form a tight, distinct cluster
        ($\sigma^2_{\rm overlap} \ll \sigma^2_{\rm noise}$) separated from noise
        in encoded space; $K \geq 2$.

  \item \textbf{Similarity:} The future result $R_{q+g}$ is \emph{more similar}
        to the current result centroid than any novel tuple; queries are
        monotonic refinements of the same analytical theme.

  \item \textbf{Random / Sampling:} $\rho \approx 1/2$ and the session is short
        ($Q < Q^*$); association rules are not yet reliable, so memory provides
        no signal and random selection is near-optimal.
\end{itemize}
\end{theorem}

\begin{table}[t!]
  \centering
  \caption{%
    Workload parameter regimes that favour each recommender.
    $\checkmark$ = condition required; $\times$ = not required
    (or reversed).
  }
  \label{tab:recipe}
  \footnotesize
  \setlength{\tabcolsep}{3pt}
  \begin{tabular}{l c c c c}
    \toprule
    \textbf{Rec.}
      & \makecell{$\rho \!\ll\! 1$\\[-2pt]\scriptsize sparse sig.}
      & \makecell{$\xi \!\gg\! 1$\\[-2pt]\scriptsize freq.\ skew}
      & \makecell{$\sigma^2_{\rm n} \!\ll\! 1$\\[-2pt]\scriptsize noise hom.}
      & \makecell{$Q \!\geq\! Q^*$\\[-2pt]\scriptsize deep hist.} \\
    \midrule
    MDI        & \checkmark & \checkmark & \checkmark & \checkmark \\
    Frequency  & $\times$   & $\times$   & $\times$   & \checkmark \\
    Clustering & $\times$   & $\times$   & $\times$\textsuperscript{a}  & $\times$ \\
    Similarity & $\times$   & $\times$   & $\times$\textsuperscript{b}  & $\times$ \\
    Random     & $\times$   & $\times$   & $\times$   & $\times$   \\
    Sampling   & $\times$   & $\times$   & $\times$   & $\times$   \\
    \bottomrule
    \multicolumn{5}{l}{\scriptsize \textsuperscript{a}distinct clusters \quad \textsuperscript{b}centroid-aligned}
  \end{tabular}
\end{table}

\paragraph{Instantiation in the benchmark generator.}
Table~\ref{tab:benchmark-params} shows how the parameters of
\texttt{generate\_benchmark.py} map to the theoretical quantities.  All four
conditions in Theorem~\ref{thm:recipe} for MDI are satisfied by construction:
$\rho \approx 0.22$ for large queries; $\xi = p_{\rm common}/(1 - p_{\rm common}) = 4$
($p_{\rm common} = 0.80$); noise rows concentrate on 2 values per attribute in a
9-attribute schema so $\sigma^2_{\rm noise} \leq 0.11$; and with $Q = 30$ queries
per session (each large query contributing $\sim 150$ tuples) the expected
antecedent support count after 10 large queries is
$10 \times 150 \times \text{supp}(r^*) \approx 90 > 1/\theta_s = 12.5$,
satisfying $Q \geq Q^*$.

\begin{table}[t!]
  \centering
  \caption{%
    Mapping from \texttt{generate\_benchmark.py} parameters to theoretical
    quantities.
  }
  \label{tab:benchmark-params}
  \footnotesize
  \setlength{\tabcolsep}{3pt}
  \begin{tabular}{l l l}
    \toprule
    \textbf{Theory param.} & \textbf{Generator param.} & \textbf{Value} \\
    \midrule
    $\rho$ (signal frac.)     & rule rows / lg.\ query sz.       & $\approx 0.22$ \\
    $\xi$ (freq.\ skew)       & $p_{\rm common} = 0.80$ bias     & $4.0$ \\
    $\sigma^2_{\rm noise}$    & 2 common vals / 9 attrs           & $\leq 0.11$ \\
    $\text{supp}(r^*)$        & 12\% rule rows; conf $= 0.85$    & $\approx 0.10$ \\
    $Q^*$                     & min\_supp $= 0.08$, 3000 rows    & $\approx 3$ lg.\ queries \\
    Core overlap $m_c$        & 20 persistent rule tuples         & $20$ \\
    Sm.\ query overlap        & 24--30 of 40--50 tuples           & $\approx 0.60$--$0.70$ \\
    \bottomrule
  \end{tabular}
\end{table}

\paragraph{Implications for real-world workloads.}
Real EDA sessions over operational databases frequently exhibit high $\xi$ simply
because data skew is ubiquitous, but they rarely satisfy $\rho \ll 1$ unless the
analyst is explicitly drilling into a rare sub-population.  This means MDI's
advantage is most pronounced in \emph{drill-down} scenarios --- exactly the kind
of session studied in the SDSS and SIMBA benchmarks --- and shrinks toward the
level of simpler baselines in broad exploratory sweeps where overlap is spread
uniformly across result rows.  The framework thus not only explains the benchmark
results but also signals where MDI deployment will deliver the greatest practical
benefit.
